\begin{trapbox}
\textbf{易错点一：同一个形式，不同的几何意义。}
\begin{itemize}
    \item 若 $P(x_0,y_0)$ 在椭圆\textbf{上}，$x_0x/a^2+y_0y/b^2=1$ 是\textbf{$P$ 处的切线}。
    \item 若 $M(x_1,y_1)$ 在椭圆\textbf{外}，同型方程是\textbf{两个切点的连线（极线）}，不是“经过外点的切线”。外点 $M$ 一般\textbf{不在极线上}：将 $M$ 代入极线左边得到 $S>1$，不是 $1$。
    \item 若 $M$ 在椭圆\textbf{内且不为原点}，极线形式是一条与椭圆不相交的实直线；不能虚构两条实切线。若 $M=(0,0)$，式子 $0=1$ 不表示普通平面直线。
\end{itemize}

\medskip
\textbf{易错点二：斜率公式不能替代切点式。}

由 $y'= -b^2x/(a^2y)$ 得出的 $k=-b^2x_0/(a^2y_0)$，\textbf{要求 $y_0\ne0$}。在 $(\pm a,0)$，切线是 $x=\pm a$，斜率不存在；不要把分母为零理解为“没有切线”。在 $(0,\pm b)$，切线是 $y=\pm b$，斜率为零。

\medskip
\textbf{易错点三：不要把半径方向直接当成椭圆法线。}

圆的切线垂直半径，但椭圆在 $P(x_0,y_0)$ 处的法向量是
\[
 \left(\frac{x_0}{a^2},\frac{y_0}{b^2}\right),
\]
一般\textbf{不平行}于 $(x_0,y_0)$。例如在 $x^2/8+y^2/2=1$ 的 $P(2,1)$，法向量为 $(1/4,1/2)$，与 $OP=(2,1)$ 显然不平行。

\medskip
\textbf{易错点四：核对椭圆方程和坐标是否匹配。}

若分母是 $a^2,b^2$，代换后还是这两个分母，不能误写成 $a,b$，也不能把 $x_0x$ 写成 $x_0^2x$。使用\textbf{已知切点}的公式前，应先检验 $x_0^2/a^2+y_0^2/b^2=1$。

\medskip
\textbf{易错点五：切线判定式里的平方、常数和前提。}

直线 $Ax+By+C=0$ 与椭圆相切的判定为
\[
 C^2=a^2A^2+b^2B^2,\qquad(A,B)\ne(0,0).
\]
不能漏掉平方，也不能把 $a^2A^2+b^2B^2$ 写成 $A^2/a^2+B^2/b^2$。把直线改写成 $Ax+By=t$ 时应使用 $t^2=a^2A^2+b^2B^2$。

\medskip
\textbf{易错点六：名称“切点幂”不是“点幂定理”。}

本条核心是\textbf{椭圆切点式与极线}。所谓“平方换乘积”是记忆技巧；它与初中圆中的“割线段乘积等于切线长平方”的\textbf{点幂定理}不是同一命题。检索可以保留“切点幂代换”这个别名，但正式推导中优先写“切点式”“极线”“切点弦”。

\medskip
\textbf{一分钟自检：}看见一个点 $M$，先计算
\[
 S=\frac{x_M^2}{a^2}+\frac{y_M^2}{b^2}.
\]
$S=1$ 用切点式；$S>1$ 求切点弦可用极线；$0<S<1$ 无实切线；$S=0$ 时形式式退化。\textbf{先判位置，再用公式。}
\end{trapbox}
